\documentclass[11pt]{amsart}
\usepackage[margin=1in]{geometry}
\usepackage{amssymb,mathtools,booktabs}
\usepackage{xurl}
\usepackage[colorlinks=true,linkcolor=blue,citecolor=blue,urlcolor=blue]{hyperref}
\newtheorem{proposition}{Proposition}

\title[Re-execution of the four-state Potts sign certificate]{A third-party re-execution of the exact sign certificate for the ferromagnetic four-state Potts reconstruction threshold}
\author{[AUTHOR NAME]}
\date{October 9, 2026}

\begin{document}
\begin{abstract}
The OpenAI Math Release preprint \emph{The Reconstruction Threshold for the Ferromagnetic Four-State
Potts Model} proves nonreconstruction at and below the Kesten--Stigum bound $d\lambda^2\le1$. Its proof
is computer assisted: via Proposition~3.1 and Lemma~6.1 it rests on two exact-arithmetic subdivision
traversals (Proposition~6.2), whose counts (Table~2) the authors obtained with GCC~13.3.0. The bundled lightweight checker validates the
recorded logs but does not repeat the traversals. We re-ran both full traversals on a different
machine with another GCC version (9.5.0) and Boost~1.86, and obtained stderr traces that are
byte-identical to the published ones. We also re-ran traversal (a) in full with Clang~17.0.6
(\texttt{-O2}) and GCC~12.3.0 (\texttt{-O1}), again byte-identical, and the first
part of traversal (b) with Clang. This is a reproducibility record. It contains no new mathematics, and it
does not check the reduction from the theorem to the coefficient arrays.
\end{abstract}
\maketitle

\noindent\textbf{Attribution.} The theorem, the reduction and the verifier
\path{verification/checks/certificate.cpp} are OpenAI's work (family 229 of
\url{https://github.com/openai/math}, commit \nolinkurl{fd4aeeb2ee4fc729c18d98444fed42fd0529eeeb},
preprint \cite{oai}). We only executed the published program and compared its output.
Developed with AI assistance (Anthropic Claude, Opus 5.5) under [AUTHOR NAME]'s direction; an
independent adversarial review was run by a separate AI session and is not human peer review.

\section{What the certificate is}
For the channel $P_\lambda(j\mid i)=\lambda\mathbf 1_{i=j}+(1-\lambda)/4$ on $[4]$, \cite[Theorem~1.1]{oai}
states that $a_n(d,\lambda)\to0$ whenever $d\lambda^2\le1$, on $d$-ary and Poisson$(d)$ trees. The
probabilistic argument reduces this to two polynomial inequalities on a product of simplices
\cite[Proposition~3.1]{oai}. These are certified by Bernstein-type coefficient signs on a midpoint
subdivision \cite[Lemma~6.1, Proposition~6.2]{oai}. The verifier is a single C++17 file (185 lines,
exact integer and \texttt{boost::multiprecision::cpp\_int} arithmetic, \texttt{assert} enabled). It
prints one progress line every 100 visited cells (mode, cell count, current depth, Type~P and Type~T
leaves so far) and finally
\begin{center}
\texttt{done 4479 19 lp=2074 lt=166}\quad (row a),\qquad \texttt{done 1935 20 lp=871 lt=97}\quad (row b),
\end{center}
which are the counts of \cite[Table~2]{oai}. The trace samples cumulative counts every 100 cells of a
deterministic depth-first traversal, so equality of full traces is strong evidence (though not a proof)
that the same cells were visited and accepted in the same order.

\section{What was already on record}
The upstream file \texttt{checks/full-runs.json} records the authors' runs: GCC~13.3.0 (Ubuntu
13.3.0-6ubuntu2\textasciitilde24.04.1), flags \texttt{-std=c++17 -O3}, elapsed $98.1$\,s and
$1287.9$\,s, with log hashes. The bundle also contains \path{mode-a-ubsan.log}, identical to
\path{mode-a.log} and not referenced anywhere, apparently a second run of (a) under UBSan. The lightweight checker \texttt{checks/verify.py} rebuilds small parts
under UBSan, checks the source binding and the recorded logs, and prints \emph{``Full traversals were
not repeated''}. So the counts rest on the authors' own executions. We found no third-party record
of a full traversal; a community project (\texttt{cert-machine}) lists this preprint as not yet checked
and plans a clean-room re-implementation.

\section{What we ran}
Machine: Intel Xeon Gold 6230R, Linux 6.8 (Ubuntu 22.04.5), single thread per run, Boost~1.86
headers. Source: \texttt{certificate.cpp}, SHA-256
\begin{center}\texttt{9f9765fb4dee9f2fc618390df89d3abc\,f2602569709fd4c2133a201b6e86bff}\end{center}
(equal to the hash in \texttt{full-runs.json}). All builds use \texttt{-std=c++17 -UNDEBUG}.
\begin{center}
\small
\begin{tabular}{lllll}
\toprule
Date & Compiler & Opt. & Run & Result (user time)\\
\midrule
2026-10-07 & GCC 9.5.0 & \texttt{-O3} & (a) full & trace identical (5m51s)\\
2026-10-07 & GCC 9.5.0 & \texttt{-O3} & (b) full & trace identical (45m06s)\\
2026-10-09 & Clang 17.0.6 & \texttt{-O2} & (a) full & trace identical (2m54s)\\
2026-10-09 & GCC 12.3.0 & \texttt{-O1} & (a) full & trace identical (4m00s)\\
2026-10-09 & Clang 17.0.6 & \texttt{-O2} & (b) 30 min & 8-line prefix identical (19m40s)\\
\bottomrule
\end{tabular}
\end{center}
``Trace identical'' means that the complete standard-error output is byte-for-byte equal to the
upstream \texttt{checks/logs/mode-a.log} (resp.\ \texttt{mode-b.log}). These files are equal to
\texttt{paper-mode-a.log} (resp.\ \texttt{paper-mode-b.log}) with the authors' trailing
\texttt{time} output removed. The 2026-10-07 binary was identified afterwards: rebuilding with
\path{g++ -std=c++17 -O3 -UNDEBUG checks/certificate.cpp} in the same directory reproduces its
SHA-256 exactly. The helper \texttt{small\_polynomials.py} (the four scalar minima of \cite[Lemma~5.1]{oai} and the
24 coordinate permutations of \cite[Section~5]{oai}) gives JSON equal to the recorded file
\path{paper-small-polynomials.json}.

The full traversal (b) was not repeated with a second compiler, because it takes more than 30 minutes
on this machine. The script \texttt{scripts/run\_all.sh} reruns the hash checks, the trace comparisons,
a fresh full traversal (a) with a compiler of the user's choice, and optionally (b) in full or as a
time-limited prefix.

\section{Scope}
This record shows that the published program, run as published, produces the published counts on
an unrelated toolchain and machine, and at three optimization levels. It does not check that the
program computes what \cite[Sections~3--6 and Appendix~A]{oai} say it computes. It also does not
check the reduction from Theorem~1.1 to the coefficient arrays. Nor is it an independent
re-implementation of the subdivision. Because the arithmetic is exact, a compiler-independent trace
is the expected outcome. Its value is evidence against undefined behaviour that changes results on the
toolchains tested, and against errors in the published logs. The archived traces contain no host or
compiler data; our local file times and \texttt{time} output support real runs, but a reader can
re-check only by re-running (mode (a) takes a few CPU minutes, mode (b) 20--45).

\begin{thebibliography}{9}
\bibitem{oai} OpenAI, \emph{The Reconstruction Threshold for the Ferromagnetic Four-State Potts Model},
OpenAI Math Release preprint, October 5, 2026,
\url{https://github.com/openai/math/tree/main/preprints/The-Reconstruction-Threshold-for-the-Ferromagnetic-Four-State-Potts-Model-October-5-2026}.
\end{thebibliography}
\end{document}
